\chapter{Number from a Fact}
\label{chap:number-from-a-fact}

The construction this book follows begins with almost nothing. It does not help itself to the
integers, or the real line, or even a primitive notion of counting. It assumes one kind of object,
and it builds everything that comes after --- arithmetic, an order on the numbers, later a calculus
of variations and a small finite field theory --- by composing copies of that one object. The
discipline of the whole development is that nothing may be introduced before it can be produced. So
the right place to start is with the single object, the first number it yields, and the one
operation without which none of the rest can even begin: telling two things apart.

This chapter does three things. It defines the primitive object, which we will call a Fact. It shows
the number that a Fact makes, and the order that the numbers carry. And it isolates distinguishing
--- the act of certifying that one carrier is not another --- as a separate, decidable capability
rather than something assumed for free. Along the way we draw on three historical experiments held in
the development's own library, each a finite and self-contained vignette: Berkeley and Galileo on
what may be introduced into a theory at all, Hume on what it means for a claim to be true, and Euclid
on why a distinction, once made, cannot be unmade.

\section{The Fact}
\label{se:the-fact}

A \emph{Fact} is a proposition together with a procedure that decides it. Both halves are essential,
and the second is the unusual one. An ordinary mathematical proposition is a statement that is either
true or false; whether anyone can settle which is a separate matter, often a hard one, sometimes an
impossible one. The construction refuses that separation. It never works with a bare proposition. It
works only with a proposition that arrives carrying the means of checking itself --- a claim and, in
the same breath, a finite computation that returns the claim's truth-value on demand.

The simplest such object asserts that true equals true. As a proposition this is trivial; as a Fact
it is the cornerstone. Its decision procedure returns ``true'' at once, by reflexivity --- the bare
observation that a thing equals itself. There is nothing to compute and nothing to wait for. What
matters is not the difficulty of the claim but its shape: a proposition and a verdict, fused. We will
treat that verdict as the reading of a real instrument. The instrument is whatever finite process
actually carries out the decision --- a physical decider, of the kind the development's experimental
library models again and again --- and asking ``is true equal to true?'' is putting a claim to that
instrument the way a voltage is put to a meter. The answer comes back, and it comes back at a cost
--- here, the smallest cost there is.

It is worth dwelling on why the decision procedure must be welded on, rather than left to be supplied
later if and when convenient. The answer is a constraint on what may legitimately enter a theory at
all, and it is old.

The development records this constraint as a single finite experiment, which it files under the
joint names of Berkeley and Galileo. Berkeley's objection to the calculus of his day was that it
trafficked in quantities --- infinitesimals, fluxions --- that no finite procedure could ever
produce, manipulate, or exhibit. They were invoked, used, and then made to vanish, and Berkeley's
complaint was that a quantity you can neither display nor distinguish is not a quantity at all, only
a convenient ghost. Galileo had pressed the complementary point from the other side: an admissible
claim about nature must be tied to an operation that leaves a recoverable trace, so that the claim
can be checked by repeating the operation. Put together, the two make one rule. \emph{Structure may
not be introduced into a theory faster than measurement can recover it.}

Read against the Fact, the rule is exactly the demand that the decision procedure travel with the
proposition. A Fact is a piece of structure that cannot outrun its own recovery, because the
proposition and the procedure that settles it are not two things but one. There is no moment at which
the claim exists but the means of checking it does not. The experiment makes the limiting case vivid:
below the resolution of a given instrument, two different descriptions of a system produce identical
records, and so, for that instrument, they are the same Fact. A distinction that no procedure can
draw is not a fact about the world; it is an artifact of a description that was introduced too fast.
The construction simply declines to ever introduce one. Everything it builds will be assembled from
objects that carry their own verification, and the chapters ahead can be read as a long exercise in
respecting that single restriction without exception.

\section{The number a Fact makes}
\label{se:the-number-a-fact-makes}

Once we have Facts we can make numbers, and we make them in the most parsimonious way available: a
number is a finite string of binary digits, and each digit is itself a Fact. There is a terminal
digit, which carries a single Fact and closes the string, and there is a prepending step, which puts
one more Fact-bearing digit in front of an existing number. Counting, in this scheme, is not a
primitive operation handed down from on high; it is the repeated act of laying one more
Fact-bearing digit alongside the ones already laid. The numerals are written big-endian, most
significant digit first, because the comparisons we care about resolve from the high end and we would
like to stop reading as early as the answer allows.

This way of building a number deserves a name for the step that does the work. The terminal digit is
a floor; every larger number is reached by the prepending step, which lays exactly one more
Fact-bearing digit in front of what is already there. That step is succession, and the development
takes it as the primitive act --- not addition, not the real line, but the bare laying-down of one
more unit. The construction files this stance under the joint names of Peano and Kushim. Peano
grounded arithmetic by making counting primitive, assuming the natural numbers into existence rather
than deriving them from anything prior. Kushim is older and more literal: the earliest personal name
that survives anywhere belongs not to a hero in a story but to a clerk, recorded on a clay tablet
tallying receivables --- a name that enters history as the observer writing to a ledger. The two mark
one principle, and the development adopts it without apology: a thing is admitted to exist exactly
when it increments the count. Each increment is a single unit and is irreversible; there are no
fractional steps, no negative steps, no quiet cancellations. An outcome that cannot be reached by
laying down one more whole digit is not a number in this system, because it is not something the
ledger can record. Existence here is not granted by explanation; it is granted by being counted.

The genuinely unusual move is what the digit's Fact is \emph{for}. It is not decoration. The
truth-value that the digit's decision procedure returns is read as the digit's sign. A Fact that
decides ``true'' contributes one way; a Fact that decides ``false'' contributes the other. So the
numbers are signed from birth, and the sign is not an extra field bolted onto a magnitude --- it is
the verdict of the very procedure the digit was built around. The order relation on these numbers is
then defined by walking the two numerals together and consulting, digit by digit, what the carried
Facts decide. A number that terminates immediately sits at or below everything. A number that
continues sits above a number that has already terminated. And when two numbers both continue, the
comparison branches on the two leading signs: agreeing positive signs pass the question down to the
remaining digits unchanged; agreeing negative signs pass it down reversed, because among negatives
the larger magnitude is the smaller number; and disagreeing signs settle the matter on the spot.

This is an ordinary lexicographic comparison, but with the comparison key supplied at each step by a
decision procedure rather than read off a static table, and it has a consequence worth naming. The
underlying two-valued logic carries only two truth-symbols, true and false. The numbers, because they
compare \emph{pairs} of carried Facts, effectively distinguish three situations: both digits true,
both digits false, and the two digits disagreeing. The first two are the cases where the signs match
and the magnitudes carry the comparison; the third is where the signs themselves decide it. One can
think of the matched cases as a single ``same'' cadence and the mismatched case as a ``different''
cadence --- a clock-complement reading in which what is recorded is not the absolute truth of each
digit but whether the two digits agree. The number system thus quietly carries one more distinction
than the two-valued logic it is built on, and that extra distinction --- agreement versus
disagreement, riding on top of true versus false --- is the seed from which several later structures
grow.

A small example makes the comparison concrete. Take the numeral whose digits all decide ``true,'' a
short run of positive digits standing for a modest positive magnitude, and compare it against a
second numeral agreeing with it on every digit but carrying one more positive digit besides. Walking
the two big-endian, the procedure finds agreement at the high end and is passed downward unchanged,
exactly as the matched-positive case prescribes, until it reaches the place where one numeral has
terminated and the other still continues; there the continuing number is ranked above the one that
stopped, and the walk halts with a verdict. Now flip the leading signs of both numerals to ``false.''
The same walk runs, but the matched-negative rule reverses the question at each step, and the number
that was the larger as a positive magnitude is now reported as the smaller --- which is exactly what
we want, since the more negative of two numbers is the lesser. Nothing in this required a stored table
of values or a notion of subtraction. It required only that each digit, when asked, returns a verdict,
and that the procedure knows how to fold two verdicts together. The order is computed, digit by
decided digit, from the high end down.

What does it mean to call any of this true? The development is careful here, and it borrows its
caution from Hume. Hume's argument, recorded as another finite experiment, is that no finite run of
observations can logically force a universal claim. A rule that has held for every recorded instance
so far is not thereby guaranteed to hold at the next instance; confirmation accumulates evidence but
never converts into necessity, because nothing in the record of the past compels the future to
resemble it. Translated into the development's vocabulary, the ledger of decided Facts grows only by
recording specific verdicts, never by recording a general law, and a rule extracted from the ledger
at one stage may not be imposed as a constraint on the next stage. Truth, for a finite observer, is
therefore not a possession but a standing: a claim earns its standing by surviving attempts to break
it under refinement, and ``certainty'' is replaced by persistence, the amount of history a claim has
so far withstood. This is precisely why a Fact must \emph{decide} its proposition rather than merely
assert it. The verdict is computed afresh whenever it is needed; it is a measurement taken now, not a
guarantee inherited from before. The number a Fact makes is, in the same spirit, a record of verdicts
and not a Platonic quantity, and the order on the numbers is retrospective --- a report on the Facts
in hand, never a decree about Facts not yet recorded.

\section{Telling carriers apart}
\label{se:telling-carriers-apart}

A number is a value, but a value has to be carried by something, and the development is explicit about
the carrier. It pairs each value with what we will call a carrier process: a Fact serving as the
value's symbol, the value itself, and a step that advances the value to its successor. The carrier
process is the development's sensor for what counts as a type of carrier in the first place. It does
not start, as the textbook natural numbers do, from the empty set; it starts from the trivial truth
that true equals true, and it stacks Facts on top of that floor. The point of isolating the carrier
is that we now have two things in hand that are easy to confuse and must not be: the value, which is
an abstract number, and the carrier, which is the concrete thing on which that value is written. The
whole of measurement will turn on keeping them apart.

The step that advances the value is worth pausing on, because it is the carrier's only moving part.
Given the current value it produces the next one, and it does so by the same logic that built the
numbers to begin with: a terminal value advances to the trivial truth that floors the system, and a
continued value advances by carrying its Fact forward onto the value beneath it. This is the ordinary
fencepost of counting --- the successor of a thing is computed from the thing, one digit at a time ---
and it is why the carrier can generate an entire run of values from a single starting Fact without
ever being handed the numbers from outside. But generating values is not the same as telling them
apart, and the development is scrupulous about the difference. The value is what is written; the
carrier is what it is written on; and a measurement is only ever a reading of a carrier, never a
direct grasp of a value. Conflating the two is the original sin the rest of the book is organized to
avoid, because a surprising amount of the confusion downstream --- about superposition, about
identity, about what a particle is --- turns out on inspection to be a value quietly mistaken for its
carrier, or a carrier mistaken for its value.

Keeping them apart requires an operation, and it is the operation this section is named for. Before
you can count occurrences of a thing, you must be able to certify that two presentations are, or are
not, the same thing. The development gives this its own capability, which it calls being
distinguishable. A family of carriers is \emph{distinguishable} when there is a decidable predicate
that separates them: a fixed symbol, a predicate that holds of any candidate differing from that
symbol, and --- the load-bearing requirement --- a guarantee that this predicate is decidable, so
that the procedure will actually return, for any candidate put to it, whether or not it differs. The
mechanism the development uses to realize this is a level structure on the carriers, a refinement
ordering of the kind the experiments' shared scaffolding records as it decodes finer events into
coarser ones: carriers that sit at different levels are, by construction, not the same, and the
difference can be read off from the levels alone, without any appeal to the carriers' contents.
Distinguishing, in other words, is not assumed; it is a procedure, and like every other capability in
the construction it is required to terminate with a verdict.

These two acts deserve their names here, at their first appearance, because the whole book is built from
them. To certify that two carriers differ --- to individuate one from another by the decidable difference
--- is to \emph{tange} them: to select a carrier out by a characteristic a procedure can actually check.
Its complement, gathering the presentations that the same predicate cannot tell apart so they may be
counted as occurrences of one thing, is to \emph{funge} them: to pool by a shared characteristic. Tange
and funge --- telling apart and pooling together, the two faces of the one decidable predicate --- are the
elementary acts every later measurement is assembled from, and this section is where the construction
first performs them.

It is tempting to treat distinguishing as too cheap to mention, but the development insists on the
opposite, and it does so through Euclid. The experiment filed under Euclid's name concerns
permanence. Euclid's constructions proceed by introducing relations --- a point, a line, an
incidence --- that, once drawn, remain available to every later step and cannot be erased without
contradicting what was built on them. The development reads this as a property of its own record. A
refinement excludes some previously open possibilities; because refinements cannot be undone, every
later observation is bound to respect every distinction already recorded. No extension of the ledger
may negate, reverse, or quietly forget a distinction it earlier drew. And from that irreversibility
comes something that looks, from a distance, like an object: what we experience as an enduring thing
is not a primitive of the world but the invariance of a recorded distinction across all the
refinements that follow it. The decidable difference of this section is the atom of that permanence.
A single certified ``these two are not the same'' is a distinction entered into the record, and the
construction's rule that the record may grow but never retract is what lets such atoms accumulate
into the stable carriers --- and, much later, the stable particles --- that the development claims to
find.

\section{What is demonstrated, and what is assumed}
\label{se:demonstrated-assumed-ch1}

Because this book takes the construction seriously as mathematics, it owes the reader, at the end of
each chapter, a plain separation of what was built and proved from what was named and modeled. The
two read alike on the page and must never be conflated, so we state the separation directly. It is
the same discipline the experimental library applies to itself: every experiment there is a finite
model that declares its own ceiling --- how much it is entitled to claim --- and is candid about what
it does \emph{not} claim. This book adopts that line chapter by chapter.

What is demonstrated in this chapter is modest and solid. There is an object, the Fact, that bundles
a proposition with a procedure deciding it, and there is a concrete simplest instance of it whose
verdict is returned, immediately, by reflexivity. There is a type of numbers built entirely from
Facts, together with an order relation defined by consulting the carried decision procedures; the
relation is a total, computable comparison, and it does on these Fact-built numerals exactly what a
signed big-endian comparison should. And there is a decidable notion of difference, with the
decidability discharged rather than assumed, so that the procedure genuinely answers when asked
whether two carriers differ. None of these requires anything beyond the construction itself: they are
definitions and the immediate consequences of definitions, established by the same finite decision
procedures whose verdicts they concern.

It is worth seeing the three results as one. A Fact decides its proposition; two carriers are told
apart by a decidable predicate; the order on the numbers is computed by consulting decisions digit by
digit. Decidability is not three separate conveniences here but a single thread, and it is the thread
the whole construction hangs from. Wherever the development later claims to have built something ---
a rational, a limit, a field --- the claim will in the end cash out as the existence of a procedure
that returns a verdict in finite time. The reason this book can speak of \emph{measuring} a
construction, rather than merely admiring it, is that every object in it is at bottom a finite
procedure for producing decisions, and a procedure that produces decisions is a thing one can put a
meter on. That a verdict also has a cost --- that the trivial Fact is settled for almost nothing
while later propositions will exhaust real budget before they resolve --- is what turns the meter from
a metaphor into a measurement, and the chapters ahead spend that budget carefully.

\begin{quote}
\emph{A note on the labels.} The physical and philosophical names in this chapter --- reading a
digit's decided truth as a \emph{sign}, calling a decision procedure's verdict a \emph{measurement},
treating an invariant recorded distinction as the germ of an \emph{enduring object} --- are
interpretations laid over the formal objects, not theorems about them. The construction earns the
mathematics: the Fact, the numbers, the order, the decidable difference. The bridge from those to the
vocabulary of measurement is a modeling choice, and it carries exactly the weight of the analogy that
licenses it, no more. Where later chapters attach heavier names --- an electron, a field, a quantity
of radiation --- the same discipline will apply, and we will mark the seam each time.
\end{quote}

What is assumed, then, is small but real, and it is honest to name it now because everything later
inherits it. We assume that carrying out a decision procedure is a fair stand-in for a measurement
--- that to put the claim that true equals true to a finite decider, and to spend what it costs to get
the answer, is to perform an observation rather than merely to contemplate one. And we assume that
reading a decided truth-value as a sign, and an agreement-or-disagreement of two such values as a
distinction, is a faithful way to let logic carry arithmetic. These are the bridges. They are not
forced on us by the symbols; they are the choices that make the symbols mean what the title of the
book says they mean. With the primitive object in hand, the number it makes, and the ability to tell
carriers apart, the construction has what it needs to climb. The next chapter takes the single step
from these foundations up through the rest of the number tower, and on toward the continuum that
finite records are only ever allowed to approach.
